\documentclass[11pt]{article}
\usepackage[a4paper,margin=27mm]{geometry}
\usepackage{amsmath,amssymb,amsthm,xcolor,booktabs,array}
\newcommand{\PROVED}{\textcolor{green!40!black}{\textbf{PROVED}}}
\newcommand{\OPEN}{\textcolor{red!70!black}{\textbf{OPEN}}}
\newcommand{\AUDIT}{\textcolor{orange!70!black}{\textbf{AUDIT}}}
\newtheorem{theorem}{Theorem}
\newtheorem{corollary}{Corollary}
\title{Fourier--Poisson duality as a Krein, rather than Hilbert, pairing (v108)}
\date{13 August 2026}

\begin{document}
\maketitle

\section{The invariant block equation}

Let $E_\rho\oplus E_{1-\overline\rho}$ be a simple functional pair in the
odd Poisson quotient and let
\begin{equation}
 A_n=\begin{pmatrix}a&0\\0&1/\overline a\end{pmatrix},
 \qquad a=n^{1/2-\rho}.                                    \tag{1}
\end{equation}
This is the centered Witt action.  Let $P=P^*$ be the Gram matrix of a
Hermitian form for which the Fourier--Poisson adjoint is the inverse Witt
action.  The required invariance is
\begin{equation}
                         A_n^*PA_n=P.                         \tag{2}
\end{equation}

\begin{theorem}[Complete solution of the invariant Gram equation ---
\PROVED]
If $\Re\rho\ne1/2$, every solution of (2) is
\begin{equation}
                         P=\begin{pmatrix}0&c\\\overline c&0\end{pmatrix}
                                                               \tag{3}
\end{equation}
for some $c\in\mathbb C$.  It is nondegenerate exactly when $c\ne0$, and
then has signature $(1,1)$.  In particular no nonzero positive definite or
positive semidefinite invariant form exists on an off-critical functional
pair.
\end{theorem}

\begin{proof}
Write $P=(p_{ij})$.  The diagonal equations in (2) are
\[
 (|a|^2-1)p_{11}=0,
 \qquad (|a|^{-2}-1)p_{22}=0.
\]
If $\Re\rho\ne1/2$, then $|a|\ne1$, hence $p_{11}=p_{22}=0$.  The
off-diagonal multiplier is
\(
 \overline a(1/\overline a)=1,
\)
so (2) imposes no further condition beyond Hermitian symmetry.  This gives
(3).  Its eigenvalues are $|c|$ and $-|c|$.
\end{proof}

\begin{corollary}[Source positivity is exactly the missing theorem ---
\PROVED]
Suppose the source-defined Fourier--Poisson pairing on the odd quotient is
nondegenerate on each functional pair and obeys the Witt adjoint identity.
If that pairing is positive semidefinite, then every spectral character
satisfies $\Re\rho=1/2$.
\end{corollary}

\begin{proof}
An off-critical pair would be nondegenerate and have signature $(1,1)$ by
Theorem 1, contradicting positive semidefiniteness.
\end{proof}

\section{What the Tate superobject does and does not supply}

The two-dimensional even object contains the characters $0$ and $1$.  It
accounts exactly for the two polar terms in the Poisson extension.  On the
primitive subspace both polar evaluations vanish, so this even block is
removed before the odd quadratic form is tested.  Therefore its positive
metric cannot alter the signature (3) of an odd functional pair.

Equivalently, the exact Tate subtraction of v106 satisfies
$\mathcal R\mathcal H_\cap=Z\mathcal H_\cap$; hence it changes the polar
extension but not the odd quotient or its invariant Gram equation.

The remaining non-circular target is consequently source-level and can be
stated without choosing spectral vectors:
\begin{equation}
 \boxed{
  \langle [g],[g]\rangle_{\mathrm{FP}}\ge0
  \quad\text{for every }[g]\in\mathcal H_-/Z\mathcal H_\cap,
 }                                                            \tag{4}
\end{equation}
where $\langle\ ,\ \rangle_{\mathrm{FP}}$ must be obtained from the
already constructed Fourier--Poisson duality and the Witt adjoint, not
defined from the spectral sign.  Once (4) is proved, Corollary 2 gives
(8c), and v99 supplies the trace-preserving Hilbert assembly and row (d).

\begin{center}
\begin{tabular}{>{\raggedright\arraybackslash}p{.58\linewidth}
                >{\centering\arraybackslash}p{.15\linewidth}
                >{\raggedright\arraybackslash}p{.17\linewidth}}
\toprule
Statement & Status & Location \\
\midrule
Invariant Gram equation on every simple functional pair & \PROVED & (1)--(3) \\
Off-critical pairs are necessarily Krein-indefinite & \PROVED & Theorem 1 \\
Tate subtraction changes the odd block signature & false & v106 and Section 2 \\
Source positivity (4) & \OPEN & the remaining mathematical theorem \\
Condition $G$ and row (d) & \OPEN & follow from (4) and v99 \\
\bottomrule
\end{tabular}
\end{center}

\end{document}
